%% parts/p02/first-order-completeness-remainder.tex --- retained additions
\noindent\textit{Reader note.} The exact inherited chapter through source line 6048
replaces the earlier elementary-substructure and L\"owenheim--Skolem overlap.
The following separately attributed historical lines 314--371 retain only
the elementary-diagram and internal--external calibrations; they do not replace
or alter the inherited proofs.

\section{Compactness applications}
\label{sec:p02-first-order-compactness-applications}
%% Depth row D-P3-013: elementary-diagram upward LS, nonstandard models,
%% and the internal/external calibration of Skolem's paradox.

For the elementary-diagram construction and these model-theoretic
calibrations, see \cite{Hodges1993,Marker2002}.

\begin{hbdefinition}{def:p02-elementary-diagram}{Parameter expansion and elementary diagram}
For a \(\Lang\)-structure \(\mathcal M\), let
\(\Lang(\mathcal M)\) add a distinct constant \(c_a\) for each
\(a\in\dom\mathcal M\), and let \(\mathcal M_{\mathcal M}\) interpret
\(c_a\) as \(a\).  The \emph{elementary diagram}
\[
 \operatorname{ElDiag}(\mathcal M)
 \defeq
 \operatorname{Th}_{\Lang(\mathcal M)}(\mathcal M_{\mathcal M})
\]
is a theory---all parameter sentences true in that expansion---not a
structure.  Unlike the parameter-free theory
\(\operatorname{Th}_{\Lang}(\mathcal M)\), it records the truth of every
formula at each named tuple; that extra information is exactly what the
elementary-extension construction uses.
\index{elementary diagram}
\index{parameter expansion}
\end{hbdefinition}

\begin{hbtheorem}{thm:p02-upward-ls-elementary-extension}{proved=annex:p2.2.20}{Upward L\"owenheim--Skolem via elementary diagrams}
Assume the Axiom of Choice.  Let \(\mathcal M\) be an infinite
\(\Lang\)-structure and
\[
 \kappa\geq
 \max\bigl(\abs{\Lang},\abs{\dom\mathcal M},\aleph_0\bigr).
\]
Then there are a \(\Lang\)-structure \(\mathcal N\) of cardinality exactly
\(\kappa\) and an elementary embedding
\(j:\mathcal M\to\mathcal N\) whose range is proper.
\end{hbtheorem}

This is the embedding form of an elementary extension
(\cref{def:p02-w02-p02-033-two-forms}).  The inclusion form, in which
\(\mathcal M\preccurlyeq\mathcal N'\) for \(\mathcal M\) itself, is
\cref{thm:p02-w02-p02-033-upward-ls-inclusion}, derived from this theorem by
\cref{thm:p02-w02-p02-033-copy-over-image}; \cref{def:p02-w02-p02-033-forms-differ}
says how the two forms differ.

\begin{hbcorollary}{thm:p02-nonstandard-arithmetic-model}{proved=annex:p2.2.21}{A countable nonstandard elementary model of arithmetic}
Assume the Axiom of Choice.  Then the standard first-order structure
\(\mathcal N_{\mathrm{ar}}=(\N,0,S,+,\cdot)\) has a countable proper
elementary extension.  Consequently a countable model satisfies every
first-order sentence true in \(\mathcal N_{\mathrm{ar}}\) while containing
elements outside the embedded standard natural numbers.
\end{hbcorollary}

The same conclusion follows directly from compactness with one new constant
(\cref{thm:p02-w02-p02-033-nonstandard-direct}, which links to the general
construction \cref{thm:nonstandard-model}).  Further compactness
constructions of nonstandard models, each with its finite satisfiability
verified, are \cref{thm:p02-w02-p02-033-divisible} (an element divisible by
every prime) and \cref{thm:p02-w02-p02-033-infinitesimal} (an elementary
extension of the real ordered field with infinitesimals); the limits of the
method are \cref{thm:p02-w02-p02-033-fs-fails}.

\begin{hbwarning}{def:p02-skolem-paradox}{Skolem's paradox is an internal--external distinction}
Assume the Axiom of Choice.  Conditionally, if an \(\Lang_{\in}\)-theory
\(T\) containing extensionality, the power-set axiom, every separation
instance and the axiom of infinity has a model, downward
L\"owenheim--Skolem supplies a model \(\mathcal M\) of \(T\) with countable
domain.  That model satisfies the existence of an \emph{uncountable} set
because \(T\) proves it: Cantor's theorem for \(\mathcal P(\omega)\) is
derivable in \(T\), so by soundness no element of \(\dom\mathcal M\) is,
in \(\mathcal M\), an injection of \(\mathcal P(\omega)^{\mathcal M}\) into
\(\omega^{\mathcal M}\) (\cref{thm:p02-w02-p02-033-cantor-internal}).
Externally, the members of \(\mathcal P(\omega)^{\mathcal M}\) form the set
\(\operatorname{ext}_{\mathcal M}(\mathcal P(\omega)^{\mathcal M})=
\{b\in\dom\mathcal M:\mathcal M\entails b\in\mathcal P(\omega)^{\mathcal M}\}\),
a subset of the countable domain and hence countable; an injection of it into
the members of \(\omega^{\mathcal M}\) exists in the ambient set theory, but
no element of the model codes it
(\cref{thm:p02-w02-p02-033-skolem-resolution}).  There is no contradiction
and no unconditional model-existence claim here; the internal statement
quantifies over elements of the model, the external one over all sets.  See
\cite{Skolem1923SetTheory,Hodges1993}.
\index{Skolem's paradox}
\end{hbwarning}
